\documentclass[12pt,a4paper]{article}

\usepackage[utf8]{inputenc}
\usepackage[T1]{fontenc}
\usepackage{lmodern}
\usepackage{amsmath,amssymb,amsthm}
\usepackage{enumitem}
\usepackage{xcolor}
\usepackage{geometry}
\geometry{a4paper, margin=2.5cm}
\usepackage{hyperref}

\hypersetup{
    colorlinks=true,
    linkcolor=blue!70!black,
    citecolor=green!50!black,
    urlcolor=blue!70!black
}

\theoremstyle{definition}
\newtheorem{definition}{Definition}[section]
\theoremstyle{plain}
\newtheorem{theorem}[definition]{Theorem}
\newtheorem{lemma}[definition]{Lemma}
\newtheorem{proposition}[definition]{Proposition}
\newtheorem{corollary}[definition]{Corollary}
\theoremstyle{remark}
\newtheorem{remark}[definition]{Remark}
\newtheorem{example}[definition]{Example}

\newcommand{\Var}{\mathrm{Var}}
\newcommand{\dist}{\mathrm{dist}}
\newcommand{\AC}{\mathrm{AC}}
\newcommand{\BV}{\mathrm{BV}}
\newcommand{\id}{\mathrm{id}}

\title{\textbf{Log-Distance Integrability for Dissipative Flows:\\
BV Selections on Fractal Attractors}}
\author{Lukas Geiger\\
\small Independent Researcher, Bernau, Germany\\
\small ORCID: \href{https://orcid.org/0009-0005-7296-1534}{0009-0005-7296-1534}}
\date{March 2026}

\begin{document}
\maketitle

\begin{abstract}
We introduce a regularity framework for nearest-point projections
onto compact attractors of dissipative semiflows in Hilbert spaces.
The classical Lipschitz projection theorem (Federer, 1959) requires
positive reach---a condition incompatible with fractal attractor
geometry.  We replace it with a \emph{trajectory log-Lipschitz}
condition (TLL), which controls the projection modulus at the
geometrically natural scale of the attractor distance rather than
the partition mesh.  Combined with a \emph{log-distance
integrability} condition (LDI), this yields bounded-variation
selections along absolutely continuous trajectories---precisely the
regularity needed for Gronwall-type arguments in conditional
regularity criteria.  The framework applies to any dissipative
semiflow with a compact attractor of finite fractal dimension,
including the 3D Navier--Stokes equations, reaction-diffusion
systems, and damped wave equations.  Two verifiable sufficient
conditions for LDI are identified: bounded variation of the
log-distance function, and exponential sublevel time-measure decay.
\end{abstract}


% ==================================================================
\section{Introduction}
% ==================================================================

Let $H$ be a separable Hilbert space and $\mathcal{A} \subset H$ a
compact set---the global attractor of a dissipative semiflow
$\{S(t)\}_{t \ge 0}$.  A fundamental question in the qualitative
theory of dissipative PDEs is the regularity of the nearest-point
projection
\[
  P(x) \;\in\; \arg\min_{a \in \mathcal{A}} \|x - a\|_H
\]
as a function of time along a trajectory $u(t) = S(t)\,u_0$.  If the
map $t \mapsto P(u(t))$ is absolutely continuous (AC), Gronwall-type
energy estimates close automatically; if it is merely measurable,
critical error terms remain uncontrolled.

The classical theory (Federer~\cite{Federer1959}) guarantees Lipschitz
continuity of the nearest-point projection on the tubular neighbourhood
of sets with \emph{positive reach}.  However, attractors of dissipative
PDEs generically have non-integer fractal dimension and reach zero
(cf.~\cite{Temam1997}, Chapter~V).  This rules out the direct
application of Federer's theorem.

The purpose of this note is to introduce a weaker regularity class---
\emph{trajectory log-Lipschitz} projections---that is compatible with
fractal geometry and still suffices for bounded-variation (BV) control
of the projected trajectory.  The key observation is that the naive
global log-Lipschitz condition
\begin{equation}\label{eq:global-LL-intro}
  \|P(x) - P(y)\|
  \;\le\;
  C\,\|x-y\|\,\bigl(1 + \log\tfrac{R}{\|x-y\|}\bigr)
\end{equation}
does \emph{not} imply BV of $P \circ u$ for AC curves~$u$: the
partition sums $\sum_i \Delta_i(1+\log(R/\Delta_i))$ diverge
logarithmically as the mesh refines (Section~\ref{sec:counterexample}).
The correct formulation replaces the increment-scale weight
$\log(R/\|x-y\|)$ with the attractor-distance weight
$\log(R/d(t))$, where $d(t) = \dist(u(t),\mathcal{A})$.  This
yields a trajectory-adapted condition whose BV sums converge
whenever a log-distance integrability condition (LDI) holds.


% ==================================================================
\section{Setup and Notation}
% ==================================================================

Throughout, $H$ denotes a separable real Hilbert space with norm
$\|\cdot\|$ and inner product $\langle\cdot,\cdot\rangle$.

\begin{definition}[Dissipative semiflow and attractor]
\label{def:semiflow}
A \emph{dissipative semiflow} on $H$ is a family
$\{S(t)\}_{t \ge 0}$ of continuous maps $S(t): H \to H$ satisfying
$S(0) = \id$, $S(t+s) = S(t) \circ S(s)$, and possessing a
\emph{compact global attractor}
$\mathcal{A} \subset H$: a compact, invariant
($S(t)\mathcal{A} = \mathcal{A}$ for all $t \ge 0$) set that
attracts every bounded subset of~$H$.
\end{definition}

\begin{definition}[Tubular neighbourhood and distance]
\label{def:tube}
For $R > 0$, the \emph{$R$-tubular neighbourhood} of $\mathcal{A}$ is
\[
  \mathcal{A}^R
  \;:=\;
  \{x \in H : \dist(x, \mathcal{A}) < R\},
\]
where $\dist(x, \mathcal{A}) := \inf_{a \in \mathcal{A}} \|x - a\|$.
For a curve $u: [0,T] \to \mathcal{A}^R$, we write
$d(t) := \dist(u(t), \mathcal{A})$.
\end{definition}

\begin{definition}[Nearest-point selection]
\label{def:selection}
A \emph{nearest-point selection} is a measurable map
$P: \mathcal{A}^R \to \mathcal{A}$ satisfying
$\|x - P(x)\| = \dist(x, \mathcal{A})$ for all
$x \in \mathcal{A}^R$, and $P|_{\mathcal{A}} = \id$.
Such a selection exists by the Kuratowski--Ryll-Nardzewski theorem
applied to the compact-valued multifunction
$x \mapsto \arg\min_{a \in \mathcal{A}} \|x-a\|$.
\end{definition}


% ==================================================================
\section{Why Global Log-Lipschitz Fails}
\label{sec:counterexample}
% ==================================================================

We first show that the global log-Lipschitz condition
\eqref{eq:global-LL-intro}, even combined with log-distance
integrability, does not directly yield BV of $P \circ u$.

\begin{proposition}[Divergence of partition sums under global LL]
\label{prop:divergence}
Let $P: \mathcal{A}^R \to \mathcal{A}$ satisfy
\eqref{eq:global-LL-intro} with constant $C > 0$.  Let
$u: [0,1] \to \mathcal{A}^R$ be Lipschitz with
$\|u'(t)\| = 1$ a.e.  For the uniform partition
$t_i = i/n$ ($i = 0, \ldots, n$), the partition sum satisfies
\[
  \sum_{i=0}^{n-1} \|P(u(t_{i+1})) - P(u(t_i))\|
  \;\le\;
  C\,(1 + \log(nR)),
\]
which diverges as $n \to \infty$.
\end{proposition}

\begin{proof}
Each increment satisfies $\Delta_i := \|u(t_{i+1}) - u(t_i)\| = 1/n$.
By \eqref{eq:global-LL-intro}:
\[
  \sum_{i=0}^{n-1} \|P(u(t_{i+1})) - P(u(t_i))\|
  \;\le\;
  C \sum_{i=0}^{n-1} \frac{1}{n}\,
  \bigl(1 + \log(nR)\bigr)
  \;=\;
  C\,(1 + \log(nR)).
\]
As $n \to \infty$, this grows as $C\log n \to \infty$.
The supremum over all partitions (which defines $\Var(P\circ u)$)
is therefore infinite.
\end{proof}

\begin{remark}
The divergence is \emph{logarithmic}, not polynomial---much slower
than for H\"older maps (where $\sum \Delta_i^\alpha$ converges for
$\alpha < 1$).  This makes global LL ``almost'' sufficient for BV,
but the logarithmic accumulation over infinitely many partition
points defeats it.  The trajectory log-Lipschitz condition avoids
this by anchoring the log weight at the attractor distance $d(t)$,
which does not depend on the partition.
\end{remark}


% ==================================================================
\section{The Trajectory Log-Lipschitz Condition}
\label{sec:TLL}
% ==================================================================

\begin{definition}[Trajectory log-Lipschitz projection (TLL)]
\label{def:TLL}
Let $u \in \AC([0,T];\, H)$ with $u(t) \in \mathcal{A}^R$ for
all~$t$, and let $P: \mathcal{A}^R \to \mathcal{A}$ be a
nearest-point selection.  We say that $P$ satisfies the
\emph{trajectory log-Lipschitz condition} (TLL) along~$u$ if there
exist constants $C_{\mathrm{TLL}} > 0$ and $h_0 > 0$ such that for
all $t \in [0,T)$ with $d(t) > 0$ and all
$0 < h < \min(h_0,\, T-t)$:
\begin{equation}\label{eq:TLL}
  \|P(u(t+h)) - P(u(t))\|
  \;\le\;
  C_{\mathrm{TLL}}\,\|u(t+h) - u(t)\|\,
  \Omega(t),
  \tag{TLL}
\end{equation}
where $\Omega(t) := 1 + \log(R/d(t))$ is the
\emph{log-distance weight}.

On the set $\{t : d(t) = 0\}$ (where $u(t) \in \mathcal{A}$),
$P(u(t)) = u(t)$ by idempotence of~$P$.  For the transition
$d(t) = 0$, $d(t+h) > 0$, the nearest-point property gives
$\|P(u(t+h)) - u(t)\| = \|P(u(t+h)) - P(u(t))\|
\le \|u(t+h) - u(t)\|$
(since $u(t) \in \mathcal{A}$ is its own nearest point and
$P$ is a contraction on $\mathcal{A}$), so no separate bound
is required.
\end{definition}

\begin{remark}[Geometric interpretation of TLL]
\label{rem:TLL-geometry}
The condition TLL asserts that the nearest-point projection,
evaluated along the trajectory~$u$, has a local Lipschitz constant
that blows up \emph{at most logarithmically} as the trajectory
approaches the attractor:
\[
  \mathrm{Lip}_{\mathrm{loc}}(P \circ u;\, t)
  \;\lesssim\;
  1 + \log\frac{R}{d(t)}.
\]
At distance $d(t) \sim R$ from the attractor, the projection is
essentially Lipschitz (the log term is $O(1)$).  As
$d(t) \to 0$, the log weight grows, reflecting the increasing
geometric complexity of the attractor boundary at fine scales.
The logarithmic rate is the \emph{critical threshold}: faster
growth (e.g., polynomial) would make LDI fail generically, while
slower growth (bounded) would imply Lipschitz projection---too
strong for fractal geometry.
\end{remark}

\begin{remark}[TLL vs global LL]
\label{rem:TLL-vs-LL}
The global log-Lipschitz condition \eqref{eq:global-LL-intro}
does \emph{not} imply TLL.  For $h \to 0$ at fixed~$t$:
\[
  \log\frac{R}{\|u(t+h)-u(t)\|}
  \;\approx\;
  \log\frac{R}{h\,|u'(t)|}
  \;=\;
  \log\frac{R}{|u'(t)|} + \log\frac{1}{h}
  \;\to\; \infty,
\]
while $\log(R/d(t))$ remains fixed.  The global condition
produces an \emph{increment-scale} log weight that diverges
with the partition mesh; TLL produces a \emph{distance-scale}
log weight that depends only on the current geometry.
This is the structural reason why TLL yields BV while
global LL does not.
\end{remark}


% ==================================================================
\section{Log-Distance Integrability}
\label{sec:LDI}
% ==================================================================

\begin{definition}[Log-distance integrability (LDI)]
\label{def:LDI}
In the setting of Definition~\ref{def:TLL}, we say that
\emph{log-distance integrability} holds if
\begin{equation}\label{eq:LDI}
  \int_0^T \|u'(t)\|\,\Omega(t)\,\mathrm{d}t
  \;=\;
  \int_0^T \|u'(t)\|\,
  \bigl(1 + \log\tfrac{R}{d(t)}\bigr)\,\mathrm{d}t
  \;<\; \infty.
  \tag{LDI}
\end{equation}
\end{definition}

\begin{remark}[Minimality of LDI]
\label{rem:LDI-minimal}
The condition LDI is positioned precisely between the known
regularity classes:
\begin{itemize}
  \item \textbf{Lipschitz projection} ($\Omega \equiv 1$):
    LDI reduces to $\int \|u'\| < \infty$, which is automatic
    for AC curves.  This is the Federer/positive-reach regime.
  \item \textbf{H\"older projection} ($\Omega(t) \sim d(t)^{-\beta}$,
    $\beta > 0$): LDI requires $\int \|u'\|\,d(t)^{-\beta} < \infty$,
    which generically fails near the attractor.  This is why
    Man\'e projectors (H\"older inverse) do not yield BV.
  \item \textbf{Log-Lipschitz projection} (TLL):
    LDI requires $\int \|u'\|\,\log(1/d) < \infty$, which is the
    borderline case---logarithmic singularities are integrable
    under mild dynamical conditions.
\end{itemize}
The log-Lipschitz/LDI pair thus occupies the \emph{critical}
regularity window for BV selections on fractal attractors.
\end{remark}


% ==================================================================
\section{Main Result: BV Chain Rule}
\label{sec:main}
% ==================================================================

\begin{theorem}[BV chain rule for log-Lipschitz projections]
\label{thm:BV-chain}
Let $\{S(t)\}_{t \ge 0}$ be a dissipative semiflow on a separable
Hilbert space~$H$ with compact global attractor
$\mathcal{A} \subset H$.  Let $u \in \AC([0,T];\, H)$ with
$u(t) \in \mathcal{A}^R$ for all~$t$, and let
$P: \mathcal{A}^R \to \mathcal{A}$ be a nearest-point selection
satisfying TLL \eqref{eq:TLL} along~$u$.  Set $v(t) := P(u(t))$.

If LDI \eqref{eq:LDI} holds, then
$v \in \BV([0,T];\, H)$ with the explicit bound
\begin{equation}\label{eq:BV-bound}
  \Var_H(v;\,[0,T])
  \;\le\;
  C_{\mathrm{TLL}}
  \int_0^T \|u'(t)\|\,\Omega(t)\,\mathrm{d}t.
\end{equation}
\end{theorem}

\begin{proof}
The proof has three steps.

\medskip
\noindent
\textbf{Step 1 (Metric derivative bound).}\;
For a.e.\ $t \in [0,T]$ with $d(t) > 0$, the \emph{metric
derivative} of $v$ at $t$ is defined as
\[
  |v'|(t)
  \;:=\;
  \lim_{h \to 0} \frac{\|v(t+h) - v(t)\|}{|h|},
\]
whenever this limit exists.  Dividing \eqref{eq:TLL} by $|h|$
and taking $h \to 0$:
\begin{equation}\label{eq:metric-deriv-bound}
  |v'|(t)
  \;\le\;
  C_{\mathrm{TLL}}\,|u'|(t)\,\Omega(t)
  \qquad \text{for a.e.\ } t \text{ with } d(t) > 0,
\end{equation}
where $|u'|(t) = \lim_{h \to 0} \|u(t+h)-u(t)\|/|h|$ exists
a.e.\ because $u$ is absolutely continuous.

On the set $\mathcal{Z} := \{t \in [0,T] : d(t) = 0\}$:
$u(t) \in \mathcal{A}$ and $v(t) = P(u(t)) = u(t)$.
For the metric derivative at such $t$, we need
$\|v(t+h) - v(t)\|/|h|$ for small $h$, which involves
$v(t+h) = P(u(t+h))$ potentially off the attractor.
The nearest-point property gives
$\|P(u(t+h)) - u(t)\| \le \|u(t+h) - u(t)\|$
(since $u(t) \in \mathcal{A}$ minimises distance to itself),
so $|v'|(t) \le |u'|(t)$.  Since $\Omega(t) \ge 1$ for all~$t$,
the bound \eqref{eq:metric-deriv-bound} holds on $\mathcal{Z}$
as well.

\medskip
\noindent
\textbf{Step 2 (BV criterion).}\;
By \cite[Theorem~3.28]{AmbrosioFuscoPallara2000}, a continuous
function $v: [0,T] \to H$ satisfying
\[
  |v'|(t) \;\le\; g(t) \quad \text{a.e.}
\]
for some $g \in L^1(0,T)$ is of bounded variation, with
$\Var(v;\,[0,T]) \le \int_0^T g(t)\,\mathrm{d}t$.

We apply this with
$g(t) := C_{\mathrm{TLL}}\,|u'|(t)\,\Omega(t)$.
By LDI \eqref{eq:LDI}, $g \in L^1(0,T)$.  The continuity of~$v$
follows from the continuity of~$u$ (which is AC, hence continuous)
and the fact that $P$ is continuous on $\mathcal{A}^R$ (as a
nearest-point projection on a compact set in a Hilbert space;
cf.~\cite[\S4.8]{BeerBook1993}).

Therefore $v \in \BV([0,T];\,H)$ and
\[
  \Var_H(v;\,[0,T])
  \;\le\;
  \int_0^T g(t)\,\mathrm{d}t
  \;=\;
  C_{\mathrm{TLL}}
  \int_0^T |u'|(t)\,\Omega(t)\,\mathrm{d}t,
\]
which is \eqref{eq:BV-bound}.

\medskip
\noindent
The bound \eqref{eq:BV-bound} follows from Steps~1 and~2 alone.
\end{proof}

\begin{remark}[Partition-sum verification]
One can also verify \eqref{eq:BV-bound} directly via partition sums.
For any partition $0 = t_0 < \cdots < t_n = T$ with mesh $< h_0$,
TLL gives
$\sum_i \|v(t_{i+1}) - v(t_i)\|
\le C_{\mathrm{TLL}} \sum_i \Omega(t_i)
\int_{t_i}^{t_{i+1}} \|u'(t)\|\,\mathrm{d}t$.
On subintervals where $d > 0$, $\Omega$ is continuous and
$\Omega(t_i) \le \sup_{[t_i,t_{i+1}]}\Omega$; on subintervals
touching $\{d = 0\}$, the nearest-point contraction
$\|v(t_{i+1})-v(t_i)\| \le \|u(t_{i+1})-u(t_i)\| \le
\int_{t_i}^{t_{i+1}} \|u'\|\,\mathrm{d}t$
controls the contribution independently of $\Omega$.  In both
cases the partial sums are bounded by
$C_{\mathrm{TLL}} \int_0^T \|u'\|\,\Omega\,\mathrm{d}t$,
confirming \eqref{eq:BV-bound}.
\end{remark}


% ==================================================================
\section{Sufficient Conditions for LDI}
\label{sec:sufficient}
% ==================================================================

The condition LDI \eqref{eq:LDI} is a single integrability
requirement.  We provide two independent sufficient conditions,
each verifiable from dynamical information about the semiflow.

% ------------------------------------------------------------------
\subsection{Variant A: Log-distance has bounded variation}
% ------------------------------------------------------------------

\begin{proposition}[BV log-distance implies LDI]
\label{prop:variant-A}
If $\log(R/d(t)) \in \BV([0,T])$, then LDI holds.
\end{proposition}

\begin{proof}
Write $\ell(t) := 1 + \log(R/d(t))$.  Since $u' \in L^1(0,T;\,H)$
(as $u$ is AC) and $\ell \in \BV([0,T]) \subset
L^\infty([0,T])$, the product
$\|u'(t)\|\,\ell(t)$ is locally integrable:
\[
  \int_0^T \|u'(t)\|\,\ell(t)\,\mathrm{d}t
  \;\le\;
  \|\ell\|_{L^\infty(0,T)}\,\int_0^T \|u'(t)\|\,\mathrm{d}t
  \;<\; \infty.  \qedhere
\]
\end{proof}

\begin{remark}[When is log-distance BV?]
\label{rem:log-BV}
The condition $\log(R/d) \in \BV$ requires that the trajectory
does not oscillate infinitely often between different distance
scales from the attractor.  Formally, it excludes infinite
``zickzack frequency'' in $d(t)$.  For monotonically attracting
trajectories ($d(t)$ eventually decreasing), this is automatic.
For trajectories with transient oscillations, BV of $\log d$
follows if the total number of scale crossings
$\{t : d(t) = \varepsilon\}$ is finite for each
$\varepsilon > 0$---a generic property of dissipative semiflows
with isolated equilibria.
\end{remark}

% ------------------------------------------------------------------
\subsection{Variant B: Sublevel time-measure control}
% ------------------------------------------------------------------

\begin{definition}[Sublevel-time control (STC)]
\label{def:STC}
The trajectory~$u$ satisfies \emph{sublevel-time control} with
exponent $\alpha > 0$ if there exists $C_* > 0$ such that
\begin{equation}\label{eq:STC}
  \bigl|\{t \in [0,T] : d(t) \le \varepsilon\}\bigr|
  \;\le\;
  C_*\,\varepsilon^\alpha
  \qquad \text{for all } \varepsilon \in (0, R].
  \tag{STC}
\end{equation}
\end{definition}

\begin{proposition}[STC implies LDI]
\label{prop:variant-B}
If STC \eqref{eq:STC} holds with some $\alpha > 0$ and
$u' \in L^2(0,T;\,H)$, then
$\log(R/d) \in L^p(0,T)$ for every $p < \infty$, and LDI holds.
\end{proposition}

\begin{proof}
\textbf{Step 1 (Tonelli representation).}\;
Using the layer-cake identity
$\log(R/d(t)) = \int_{d(t)}^{R} s^{-1}\,\mathrm{d}s$ and Tonelli:
\begin{equation}\label{eq:tonelli}
  \int_0^T \|u'(t)\|\,\log\frac{R}{d(t)}\,\mathrm{d}t
  \;=\;
  \int_0^R \frac{1}{s}
  \Bigl(\int_{\{d(t) \le s\}} \|u'(t)\|\,\mathrm{d}t\Bigr)
  \,\mathrm{d}s.
\end{equation}

\medskip
\noindent
\textbf{Step 2 (H\"older on inner integral).}\;
By H\"older's inequality:
\[
  \int_{\{d(t)\le s\}} \|u'(t)\|\,\mathrm{d}t
  \;\le\;
  \Bigl(\int_0^T \|u'(t)\|^2\,\mathrm{d}t\Bigr)^{1/2}
  \cdot
  \bigl|\{d \le s\}\bigr|^{1/2}
  \;\le\;
  \|u'\|_{L^2}\,C_*^{1/2}\,s^{\alpha/2},
\]
where the last step uses STC.

\medskip
\noindent
\textbf{Step 3 (Integrability near $s=0$).}\;
Substituting into \eqref{eq:tonelli}:
\[
  \int_0^R \frac{1}{s}\,\|u'\|_{L^2}\,C_*^{1/2}\,s^{\alpha/2}
  \,\mathrm{d}s
  \;=\;
  \|u'\|_{L^2}\,C_*^{1/2}
  \int_0^R s^{\alpha/2 - 1}\,\mathrm{d}s
  \;=\;
  \|u'\|_{L^2}\,C_*^{1/2}\,
  \frac{2}{\alpha}\,R^{\alpha/2}
  \;<\; \infty.
\]
The integral converges because $\alpha > 0$, so
$s^{\alpha/2-1}$ is integrable near $s=0$.

\medskip
\noindent
\textbf{Step 4 ($L^p$ bounds).}\;
The same Tonelli argument with H\"older exponents $p$ and $q$
($1/p + 1/q = 1$) gives
$\|\log(R/d)\|_{L^p} \le C(p,\alpha,R,C_*)$ for every
$p < \infty$, so LDI follows from any finite $L^q$ bound
on~$\|u'\|$.
\end{proof}

\begin{remark}[Sources of STC in dissipative flows]
\label{rem:STC-sources}
Sublevel-time control \eqref{eq:STC} is the genuine new
analytical input.  Three mechanisms can generate it:

\begin{enumerate}[label=(\roman*)]
  \item \textbf{Exponential tracking:}\;
    If $d(t) \le C_0\,e^{-\lambda t}$ for $t \ge t_0$, then
    $\{d \le \varepsilon\} \subset
    \{t \ge (\log(C_0/\varepsilon))/\lambda\}$,
    whose measure is $O(\log(1/\varepsilon))$---much smaller than
    $C_*\,\varepsilon^\alpha$ for any $\alpha > 0$.
    So exponential tracking implies STC with \emph{any} exponent.
  \item \textbf{Lyapunov dissipation:}\;
    If a Lyapunov functional $E$ satisfies
    $-\dot{E} \ge \Psi(d)$ with $\Psi$ not too flat near $d=0$
    (e.g.\ $\Psi(d) \ge c\,d^\beta$ for some $\beta$),
    then sublevel times are controlled by the total energy budget
    $E(0) - \inf E$.
  \item \textbf{No fast oscillation:}\;
    Energetic constraints (finite dissipation integral)
    prevent arbitrarily rapid back-and-forth motion near the
    attractor, bounding the number of excursions into
    $\{d \le \varepsilon\}$.
\end{enumerate}
Each mechanism yields STC with different exponents~$\alpha$.
The STC condition is physically plausible for all standard
classes of dissipative PDEs.
\end{remark}


% ==================================================================
\section{Applications to Dissipative PDEs}
\label{sec:applications}
% ==================================================================

We illustrate how the abstract framework applies to concrete
PDE settings.  In each case, the semiflow possesses a compact
global attractor, and the question is whether TLL and LDI hold.

% ------------------------------------------------------------------
\subsection{3D Navier--Stokes on $\mathbb{T}^3$}
% ------------------------------------------------------------------

\begin{example}[Navier--Stokes equations]
\label{ex:NS}
Let $H = L^2(\mathbb{T}^3;\,\mathbb{R}^3)$ (divergence-free),
$\mathcal{A}$ the global attractor (in the sense of Temam~\cite{Temam1997}) with smooth
forcing~$f$.  By Temam~\cite{Temam1997}:
\begin{itemize}
  \item $\mathcal{A}$ is compact in $H^1$ (hence in $H$).
  \item Every $a \in \mathcal{A}$ is $C^\infty$.
  \item Exponential tracking holds: $d(t) \le C_0\,e^{-\lambda t}$
    for $t \ge t_0$.
\end{itemize}

The attractor has finite fractal dimension
$d_f(\mathcal{A}) \le c_1\,G^{2/3}(1+\log G)^{1/3}$
(where $G$ is the Grashof number;
cf.~\cite{ConstantinFoiasTemam1988}), but generically has
reach~$0$---ruling out Federer's theorem.

\medskip\noindent
\textbf{Status of TLL:}\; Open.  The squeezing property of the
Navier--Stokes semiflow (cf.~\cite{FoiasManleyRosaTemam2001},
Chapter~14) provides exponential contraction of high-frequency
differences on the attractor.  Combined with Man\'e projector
estimates, this suggests that the nearest-point projection has
at most log-Lipschitz modulus along trajectories---but a rigorous
proof requires quantitative control of the projection stability
in the high-mode complement $Q_N H$.

\medskip\noindent
\textbf{Status of LDI:}\; Follows from exponential tracking
on $[t_0, T]$ for any fixed $T < \infty$
(Remark~\ref{rem:STC-sources}).  The critical question is the
blow-up regime $T \nearrow T^*$, where $\|u'\| \to \infty$.
The BV bound \eqref{eq:BV-bound} yields Assumption~G$'$ of
\cite{FST-NS2026}, which implies global regularity via the
attractor-distance argument.
\end{example}

% ------------------------------------------------------------------
\subsection{Reaction-diffusion systems}
% ------------------------------------------------------------------

\begin{example}[Reaction-diffusion on bounded domains]
\label{ex:RD}
Consider $\partial_t u = D\Delta u + f(u)$ on a bounded domain
$\Omega \subset \mathbb{R}^n$ with Dirichlet or Neumann
boundary conditions, where $D > 0$ is a diffusion matrix and
$f$ satisfies a dissipativity condition
$\langle f(u), u \rangle \le C - c|u|^2$.  The semiflow
possesses a compact global attractor $\mathcal{A} \subset L^2(\Omega)$
with finite fractal dimension~\cite{Temam1997}.

For scalar equations ($u \in \mathbb{R}$) in one space dimension
with a gradient nonlinearity $f(u) = -V'(u)$, the attractor
consists of equilibria and heteroclinic orbits, and the
nearest-point projection is Lipschitz (in fact, the attractor has
a Lipschitz-graph structure as an inertial manifold;
cf.~\cite{MalletParetSell1988}).  TLL and LDI hold trivially.

For systems in higher dimensions, the attractor may have non-integer
fractal dimension and reach zero.  The log-Lipschitz framework then
provides a path to BV selections that Federer's theorem does not.
The spectral gap condition for inertial manifolds is satisfied for
many reaction-diffusion systems~\cite{EdenFoiasNicolaenkoTemam1994},
making the attractor geometry more regular than in the
Navier--Stokes case.
\end{example}

% ------------------------------------------------------------------
\subsection{Damped wave equations}
% ------------------------------------------------------------------

\begin{example}[Damped semilinear wave equation]
\label{ex:wave}
Consider $\partial_{tt} u + \gamma\,\partial_t u = \Delta u + f(u)$
on $\Omega \subset \mathbb{R}^3$ with $\gamma > 0$ and subcritical
nonlinearity.  The semiflow on $H^1_0 \times L^2$ possesses a
compact global attractor with finite fractal dimension
(cf.~\cite{BabinVishik1992}).  The exponential tracking property
holds with rate $\lambda = \gamma/2$ (spectral gap of the damped
wave operator).

The attractor geometry is intermediate between reaction-diffusion
(where inertial manifolds often exist) and Navier--Stokes (where
they do not).  The log-Lipschitz framework provides a
\emph{uniform} approach applicable across all three settings,
with TLL as the single geometric input.
\end{example}


% ==================================================================
\section{Discussion}
\label{sec:discussion}
% ==================================================================

\subsection{The new mathematical interface}

The central contribution of this note is the identification of
the \emph{critical regularity interface} for BV selections on
fractal attractors.  The interface is characterised by a single
question:

\begin{center}
\fbox{\parbox{0.85\textwidth}{%
\textbf{Core Question.}\;
For a dissipative semiflow with compact attractor
$\mathcal{A} \subset H$, does the nearest-point projection satisfy
trajectory log-Lipschitz control \eqref{eq:TLL}?
}}
\end{center}

\noindent
This is a question about the \emph{regularity of the attractor as
seen from the trajectory}, not about the global geometry
of~$\mathcal{A}$.  The distinction is crucial: global geometric
conditions (positive reach, inertial manifold) impose structure
on the \emph{entire} attractor, while TLL requires regularity
only in the direction of the dynamical approach.

\subsection{Relation to existing regularity theories}

\begin{center}
\begin{tabular}{lllll}
\hline
\textbf{Condition} & \textbf{Projection} & \textbf{Composition}
  & \textbf{Fractal?} & \textbf{Source} \\
\hline
Positive reach $> 0$
  & Lipschitz-1
  & AC $\Rightarrow$ AC
  & No
  & Federer~\cite{Federer1959} \\
Prox-regular
  & locally Lipschitz
  & AC $\Rightarrow$ AC (local)
  & No
  & Poliquin et al.~\cite{PoliquinRockafellarThibault2000} \\
Hausdorff-BV set map
  & BV selection exists
  & ---
  & Yes
  & Chistyakov~\cite{Chistyakov2004} \\
Man\'e projector
  & H\"older${}^{-1}$
  & AC $\Rightarrow$ H\"older
  & Yes
  & Man\'e~\cite{Mane1981} \\
\textbf{TLL + LDI}
  & \textbf{log-Lipschitz}
  & \textbf{AC $\Rightarrow$ BV}
  & \textbf{Yes}
  & \textbf{This note} \\
\hline
\end{tabular}
\end{center}

\noindent
The TLL+LDI framework is the only entry that simultaneously
(i)~accommodates fractal geometry (non-integer dimension, reach~$0$),
(ii)~yields BV regularity (sufficient for Gronwall-type arguments),
and (iii)~is formulated as a trajectory-level condition
(independent of global attractor structure).

\subsection{Limitations}

\begin{itemize}
  \item \textbf{TLL is open for 3D NS.}\;
    The squeezing property and Man\'e projector estimates suggest
    log-Lipschitz stability of the projection, but a rigorous proof
    requires quantitative bounds on the high-mode projection
    complement that are not currently available.
  \item \textbf{LDI is non-trivial only under blow-up.}\;
    For smooth solutions on compact time intervals, LDI holds
    automatically.  The condition becomes critical only in the
    limit $T \nearrow T^*$ (blow-up time), where $\|u'\| \to \infty$.
  \item \textbf{BV, not AC.}\;
    The framework yields BV regularity, which suffices for
    Stieltjes-integral-based Gronwall arguments but not for
    pointwise differential inequalities.  Upgrading BV to AC
    would require Lipschitz projection (positive reach)---exactly
    the condition we are trying to avoid.
\end{itemize}


% ==================================================================
\section{Conclusion}
% ==================================================================

We have introduced a regularity framework for nearest-point
projections onto compact attractors that occupies the critical
gap between H\"older (insufficient for BV) and Lipschitz
(incompatible with fractal geometry).  The trajectory log-Lipschitz
condition (TLL) controls the projection modulus at the attractor
distance scale, and the log-distance integrability condition (LDI)
ensures that the resulting BV sums converge.

The framework reduces conditional regularity criteria for
dissipative PDEs---including the 3D Navier--Stokes
equations---to a single geometric-dynamical question: does the
nearest-point projection have at most logarithmic blowup along
dissipative trajectories?  This question is independent of the
specific PDE and depends only on the interplay between attractor
geometry and the semiflow's contraction properties.


% ===================================================================
% AI Disclosure
% ===================================================================
\subsection*{AI Disclosure}
In the preparation of this manuscript, AI-assisted tools (Claude, Anthropic)
were used in a supporting capacity, particularly for literature research,
text structuring, and formulation assistance.  The conceptual design,
scientific argumentation, and responsibility for the entire content
rest solely with the author.


\begin{thebibliography}{99}

\bibitem{AmbrosioFuscoPallara2000}
L.~Ambrosio, N.~Fusco, and D.~Pallara,
\emph{Functions of Bounded Variation and Free Discontinuity Problems},
Oxford Mathematical Monographs, Clarendon Press, Oxford, 2000.

\bibitem{BabinVishik1992}
A.~V.~Babin and M.~I.~Vishik,
\emph{Attractors of Evolution Equations},
Studies in Mathematics and its Applications, vol.~25,
North-Holland, Amsterdam, 1992.

\bibitem{BeerBook1993}
G.~Beer,
\emph{Topologies on Closed and Closed Convex Sets},
Mathematics and its Applications, vol.~268,
Kluwer Academic Publishers, 1993.

\bibitem{Chistyakov2004}
V.~V.~Chistyakov,
\emph{Selections of bounded variation},
J.~Appl.\ Anal.\ \textbf{10} (2004), no.~1, 1--82.

\bibitem{ConstantinFoiasTemam1988}
P.~Constantin, C.~Foias, and R.~Temam,
\emph{Attractors representing turbulent flows},
Mem.\ Amer.\ Math.\ Soc.\ \textbf{53} (1985), no.~314.

\bibitem{EdenFoiasNicolaenkoTemam1994}
A.~Eden, C.~Foias, B.~Nicolaenko, and R.~Temam,
\emph{Exponential Attractors for Dissipative Evolution Equations},
Research in Applied Mathematics, vol.~37,
Wiley--Masson, Paris, 1994.

\bibitem{Federer1959}
H.~Federer,
\emph{Curvature measures},
Trans.\ Amer.\ Math.\ Soc.\ \textbf{93} (1959), 418--491.

\bibitem{FoiasManleyRosaTemam2001}
C.~Foias, O.~Manley, R.~Rosa, and R.~Temam,
\emph{Navier--Stokes Equations and Turbulence},
Encyclopedia of Mathematics and its Applications, vol.~83,
Cambridge University Press, 2001.

\bibitem{FST-NS2026}
L.~Geiger,
\emph{A Thermodynamic Obstruction to Finite-Time Blow-Up in the 3D
Navier--Stokes Equations (Conditional Framework)},
FST Navier--Stokes Paper, 2026.

\bibitem{MalletParetSell1988}
J.~Mallet-Paret and G.~R. Sell,
\emph{Inertial manifolds for reaction diffusion equations in higher
space dimensions},
J.~Amer.\ Math.\ Soc.\ \textbf{1} (1988), no.~4, 805--866.

\bibitem{Mane1981}
R.~Man\'e,
\emph{On the dimension of the compact invariant sets of certain
non-linear maps},
Lecture Notes in Mathematics, vol.~898, Springer, Berlin, 1981,
pp.~230--242.

\bibitem{PoliquinRockafellarThibault2000}
R.~A.~Poliquin, R.~T.~Rockafellar, and L.~Thibault,
\emph{Local differentiability of distance functions},
Trans.\ Amer.\ Math.\ Soc.\ \textbf{352} (2000), no.~11, 5231--5249.

\bibitem{Temam1997}
R.~Temam,
\emph{Infinite-Dimensional Dynamical Systems in Mechanics and Physics},
2nd ed., Applied Mathematical Sciences, vol.~68,
Springer, New York, 1997.

\end{thebibliography}

\end{document}
